\chapter{누구를 믿을까?}

새 학기 첫날이에요.
연필 끝에 달린 지우개가 다 닳았는데 필통을 열어 보니 지우개가 없어요.
옆자리 친구에게 빌려야겠는데, 반 친구들 얼굴이 아직 낯설어요.
누구에게 부탁하면 좋을까요?

첫날에는 알 수가 없어요.
아무에게나 물어볼 수밖에 없지요.
그런데 한 달쯤 지나면 이야기가 달라져요.
서연이는 빌린 물건을 꼭 돌려준다는 것,
도윤이는 착하지만 자주 잊어버린다는 것,
지우는 누구에게나 잘 빌려준다는 것을 알게 돼요.

어떻게 알게 되었을까요?
누가 알려 준 것도 아닌데 말이에요.
답은 “일어난 일”이에요.
한 달 동안 교실에서 작은 일이 수백 번 일어났어요.
누가 누구에게 무엇을 빌려주었는지, 누가 약속을 지켰는지, 누가 누구와 짝이 되고 싶어 했는지.
여러분은 그 작은 일들을 보면서 머릿속에 저절로 지도를 그렸어요.
누구를 믿어도 되는지 알려 주는 지도예요.

이렇게 여러 사람이 한 사람에 대해 갖게 되는 믿음을 \textbf{평판}이라고 해요.
“그 친구는 믿을 만해”라는 말은 평판이 좋다는 뜻이에요.

\section{어른들의 세상에도 평판이 있어요}

어른들도 평판을 아주 많이 써요.
은행은 돈을 빌려줄 때 그 사람이 예전에 빌린 돈을 잘 갚았는지 기록을 살펴봐요.
그 기록을 숫자 하나로 줄인 것을 신용 점수라고 해요.
인터넷 쇼핑을 할 때 사람들이 별점을 보는 것도 같아요.
별 다섯 개를 받은 가게는 많은 사람이 믿을 만하다고 말해 준 가게예요.

이런 점수들에는 공통점이 있어요.
점수는 하늘에서 떨어지지 않아요.
모두 \emph{일어난 일의 기록}에서 나와요.
돈을 갚은 기록, 물건을 사고 남긴 별점, 약속을 지킨 기록에서요.

\section{얼굴이 보이지 않는 곳}

교실에서는 친구 얼굴을 볼 수 있어요.
그런데 인터넷에는 얼굴이 없어요.
특히 이 책에서 이야기할 \textbf{블록체인}이라는 곳에서는 이름조차 없어요.
사람들은 \texttt{0x7a3f...}처럼 숫자와 알파벳이 길게 이어진 이름표로만 서로를 알아봐요.
그 이름표 뒤에 어른이 있는지, 컴퓨터 프로그램이 있는지도 알 수 없어요.

그래도 한 가지는 남아 있어요.
\emph{기록}이에요.
블록체인에서는 누가 누구에게 무엇을 보냈는지, 누가 누구에게 무엇을 허락했는지가
하나도 빠짐없이 적혀 있고, 누구나 그 기록을 읽을 수 있어요.

그렇다면 이런 질문을 해 볼 수 있어요.

\begin{center}
\begin{tcolorbox}[colback=white, colframe=inkblue, boxrule=0.8pt, arc=2mm, width=0.9\linewidth, halign=center]
\titlefont\bfseries 기록만 보고, 누구를 믿을 만한지\\점수로 나타낼 수 있을까?
\end{tcolorbox}
\end{center}

이것이 제 논문의 질문이에요.
그리고 이 책은 제가 그 질문에 어떻게 답하려고 했는지,
무엇이 잘되었고 무엇이 잘 안 되었는지에 대한 이야기예요.

\begin{think}
여러분 반에서 “믿을 만한 친구”를 한 명 떠올려 보세요.
그 친구를 믿게 된 일을 세 가지만 적어 볼 수 있나요?
그 일들은 모두 “누가 누구에게 무엇을 했다”는 모양을 하고 있지 않나요?
\end{think}

\begin{learned}
\begin{itemize}
\item 평판은 여러 사람이 한 사람에 대해 갖는 믿음이에요.
\item 평판은 일어난 일의 기록에서 생겨요.
\item 블록체인에는 얼굴도 이름도 없지만, 기록은 모두 남아 있어요.
\end{itemize}
\end{learned}
